% !TEX root = ../../main-es.tex
% appendix/es/E-tipos-operadores.tex
% Longtable multipage con tipos y operadores del fragmento .os formalizado.
% Estilo: booktabs + dcolumn (alineación decimal para aridades) + endhead/endfoot.

\chapter{Tipos y operadores del fragmento \dsl{.os}}
\label{app:tipos-operadores}

\porque{Referencia rápida.}{%
  Tabla de referencia del fragmento \dsl{.os} cuyo marco formal vive en el
  cap.~\ref{cap:marco-formal}. Cada fila ata un operador a su kind, su aridad y
  la categoría denotacional que lo modela. La columna \emph{Ref.} apunta a la
  sección donde se define.}

\section{Catálogo de operadores}

La tabla~\ref{tab:ops-os} cataloga los operadores del fragmento \dsl{.os}
incluidos en la semántica denotacional (cap.~\ref{cap:marco-formal}). Los
operadores se agrupan por \emph{kind}: \textsc{Control} (plano determinista),
\textsc{Effect} (plano del LLM) y \textsc{Pure} (sin efecto). La aridad se
expresa como $n{\to}m$ donde $n$ es el número de snapshots de entrada y $m$ el
de salida; el peso $w$ es el coste relativo de invocación (estimado, ver
cap.~\ref{cap:experimentos}).

\begin{longtable}{@{}ll d{2.1} c c@{}}
\caption[Operadores del fragmento \dsl{.os} y su clasificación.]%
        {Operadores del fragmento \dsl{.os}: kind, aridad $n{\to}m$, peso $w$
         (invocaciones relativas) y sección de definición.}%
\label{tab:ops-os}\\
\toprule
\textbf{Operador} & \textbf{Kind} & \multicolumn{1}{c}{\textbf{Aridad}} &
\textbf{Peso} & \textbf{Ref.} \\
\midrule
\endfirsthead

\caption[]{(continuación)}\\
\toprule
\textbf{Operador} & \textbf{Kind} & \multicolumn{1}{c}{\textbf{Aridad}} &
\textbf{Peso} & \textbf{Ref.} \\
\midrule
\endhead

\midrule
\multicolumn{5}{r}{\footnotesize Continúa en la página siguiente.} \\
\endfoot

\bottomrule
\endlastfoot

\oskw{literal}   & \textsc{Pure}    & 0.0 & 0.01 & \S\ref{sec:semantica} \\
\oskw{formula}   & \textsc{Pure}    & 0.0 & 0.02 & \S\ref{sec:semantica} \\
\oskw{bif}       & \textsc{Control} & 1.1 & 0.05 & \S\ref{sec:semantica} \\
\oskw{call}      & \textsc{Control} & 1.1 & 0.10 & \S\ref{sec:semantica} \\
\oskw{ret}       & \textsc{Control} & 1.1 & 0.05 & \S\ref{sec:semantica} \\
\oskw{ask}       & \textsc{Effect}  & 1.1 & 1.00 & \S\ref{sec:semantica} \\
\oskw{ai}        & \textsc{Effect}  & 1.1 & 1.20 & \S\ref{sec:semantica} \\
\oskw{map}       & \textsc{Control} & 1.1 & 0.08 & \S\ref{sec:semantica} \\
\oskw{filter}    & \textsc{Control} & 1.1 & 0.08 & \S\ref{sec:semantica} \\
\oskw{reduce}    & \textsc{Control} & 1.1 & 0.09 & \S\ref{sec:semantica} \\
\oskw{foreach}   & \textsc{Control} & 1.1 & 0.10 & \S\ref{sec:semantica} \\
\oskw{fork}      & \textsc{Control} & 1.n & 0.15 & \S\ref{sec:semantica} \\
\oskw{join}      & \textsc{Control} & n.1 & 0.20 & \S\ref{sec:semantica} \\
\oskw{call-multi}& \textsc{Control} & 1.n & 0.40 & \S\ref{sec:semantica} \\
\oskw{hot-run}   & ---              & --- & ---  & fuera del fragmento \\
\oskw{reflect}   & ---              & --- & ---  & dirección futura (M3) \\
\end{longtable}

\section{Convenciones de la tabla}

\begin{itemize}
\item \textbf{Kind}: \textsc{Pure} no invoca al LLM; \textsc{Control} esparte
  del plano determinista; \textsc{Effect} cuarentena la invocación al LLM
  (\S\ref{sec:semantica}).
\item \textbf{Aridad} $n{\to}m$: $n$ snapshots de entrada, $m$ de salida.
  \texttt{1.n} indica fan-out (fork); \texttt{n.1} fan-in (join).
\item \textbf{Peso}: invocaciones relativas normalizadas a
  \oskw{ask}$=1{,}00$. Estimación empírica del cap.~\ref{cap:experimentos}.
\item Los operadores fuera del fragmento (\oskw{hot-run}, \oskw{reflect}) se
  listan para acotar el alcance; su semántica queda como trabajo futuro
  (\S\ref{sec:reflexion}).
\end{itemize}
